\section{Related Work}
\label{sec:related}

\subsection{Tactile Perception}

Driven by rapid progress in tactile sensor hardware~\cite{gelsight}, tactile perception has
developed quickly in recent years to complement vision-based perception. A variety of tasks
have been tackled with touch, generally targeting situations in which vision alone is
insufficient.
First of all, in object recognition, tactile information is used to discern the fine-grained surface properties of an object to decide its class and material type~\cite{textureuneven,glovetactile,triboelectrictactile,bagoffeatures,dotsim}.
A second line addresses cases where 3D reconstruction solely based on vision is challenging. 
Occlusion is one such case: during a grasp the gripper hides much of the object from the camera, and here tactile measurements could be used to recover the hidden geometry~\cite{neuralfeels,touchsdf}.
In addition, for glossy and transparent surfaces, since specular reflections violate the photometric assumptions that multi-view reconstruction relies on, tactile sensors could be used to obtain surface geometry measurements as the sensors are is unaffected by such factors~\cite{snaptapsplat,splattouch,touchgs}.
Beyond these failure cases, tactile sensors resolve
surface detail far below what a camera can see, which can be accumulated into highly detailed
reconstructions~\cite{gelslam}.
Finally, since tactile measurements carry a signature of the local material and surface geometry, it
can also be used to localize the sensor itself, either against an image or against a pre-built
3D map of the object~\cite{tacloc,splattouch,normalflow}.
Building on these prior works, our paper tackles a new task: predicting tactile \emph{videos} at
locations that have not been touched before.
This is complementary to the works above, since it can
augment an already collected set of touches and thereby improve downstream tactile tasks such as
object recognition or 3D reconstruction.

\subsection{Predicting Sensory Signals from Examples}

The ability to predict (or imagine) sensory signals is an important capacity for intelligent beings such as humans to act in the world.
Such predictions guide planning, which is the key component to solve complex tasks in everyday life.
Motivated by such facts, there have been large bodies of work in predicting various sensory signals (e.g., vision and audio) from a sparse set of examples.

\paragraph{Example-based image synthesis}
A large body of prior work studies example-based image synthesis. In its usual form, one is given
a pair of images related by an unknown transformation, such as a change of style or of texture,
and must synthesize, for a new target image, what its transformed version would look like. Early
approaches recover this mapping using classical image matching~\cite{imageanalogies,patchmatch}.
Modern approaches instead use diffusion models, or establish dense correspondences between the
source images using learned descriptors and
matchers~\cite{ldm,lightglue,asic}. Our task can be seen as an extension of this
setting in which the source images are the normal maps at each tactile sensor location, and the
target ``images'' are the high-dimensional tactile videos we aim to generate.

\paragraph{Example-based synthesis of audio and touch}
Inspired by the image literature, recent work has attempted example-based synthesis for other
sensory signals. For audio, several methods generate spectrograms by jointly conditioning on
video frames and reference audio~\cite{foleyanalogies,foleycrafter,audiotextureanalogy}. For
touch, recent work mainly generates a single tactile image at the moment of contact from
reference touch data. Tactile DreamFusion~\cite{tactiledreamfusion} uses image quilting to tile a
single tactile measurement across the remaining surface of an object. Tactile-Augmented Radiance
Fields~\cite{tarf} trains a diffusion model on reference touches, conditioned on nearby image
captures, to predict tactile images at new locations. ObjectFolder~\cite{objectfolder} and
ObjectFolder 2.0~\cite{objectfolder2} train an implicit neural representation, per object, that
takes pixel coordinates and a pressing depth as input and outputs the tactile image at that
location.

We extend this last strand of works from images to videos. This is substantially more challenging, since a
video requires modeling how the sensor's gel deforms over the course of a press rather than only
its state at one instant.
